\documentclass[conference]{IEEEtran}
\IEEEoverridecommandlockouts

\usepackage{cite}
\usepackage{amsmath,amssymb,amsfonts}
\usepackage{graphicx}
\usepackage{textcomp}
\usepackage{xcolor}
\usepackage{booktabs}
\usepackage{url}
\usepackage{microtype}
\usepackage{listings}

\lstset{
  basicstyle=\ttfamily\scriptsize,
  frame=single,
  breaklines=true,
}

\def\BibTeX{{\rm B\kern-.05em{\sc i\kern-.025em b}\kern-.08em
    T\kern-.1667em\lower.7ex\hbox{E}\kern-.125emX}}

\begin{document}

\title{Deterministic MCP Guardrails for Safe Power-Pathway\\Design Agents}

\author{%
\IEEEauthorblockN{Chenrui Xu}
\IEEEauthorblockA{\textit{Department of Civil and Environmental Engineering} \\
Carnegie Mellon University \\
Pittsburgh, PA, USA \\
chenruixu@cmu.edu}
}

\maketitle

\begin{abstract}
Large language model (LLM) agents are increasingly used for safety-critical
engineering decisions, yet they remain prone to arithmetic hallucination and
constraint omission.  We study whether deterministic, pre-execution guardrails
implemented via the Model Context Protocol (MCP) can reduce invalid
decisions in grid-constrained power-pathway planning.  We introduce a
22-case benchmark spanning five failure modes (infeasible phases, budget
violations, carbon violations, hallucination traps, and feasible baselines)
and four agent conditions that progressively add structure: prompt-only (B1),
self-critique (B2), tools without enforcement (B3), and MCP-guarded with
mandatory validator calls (G1).  Live evaluation with Gemini~2.5~Flash
shows that providing validator tools without enforcement (B3) reduces the
invalid-decision rate from 36.4\% to 9.1\%, while adding gateway
enforcement (G1) further reduces it to 4.5\%.  Self-critique provides no
improvement over the prompt-only baseline.  In this benchmark, tool
availability accounts for most of the improvement, with pre-execution
enforcement adding a smaller but measurable safety margin.
\end{abstract}

\begin{IEEEkeywords}
LLM agents, guardrails, model context protocol, power systems, constraint
enforcement, safe AI
\end{IEEEkeywords}

% ============================================================
\section{Introduction}
% ============================================================

Data-center operators and grid planners are beginning to adopt LLM-based
agents to recommend power-supply strategies, selecting among technologies such
as grid interconnects, wind-plus-battery, solar-plus-battery, fuel cells,
and small modular reactors~\cite{iea2024datacenters}.  These recommendations
must satisfy hard engineering constraints: firm (dispatchable) capacity must
cover peak load in every growth phase; total capital expenditure must fit
within a budget ceiling; blended carbon intensity must not exceed a regulatory
target; and system uptime must meet a reliability threshold.

Unfortunately, current LLMs are unreliable arithmetic reasoners~\cite{frieder2023mathematical},
and their self-monitoring through chain-of-thought or self-critique
is insufficient for safety-critical domains~\cite{huang2023large}.
Even when a model ``checks'' its own work, it may reinforce incorrect
reasoning rather than correct it.

The \emph{Model Context Protocol} (MCP)~\cite{anthropic2024mcp} provides a
standardized interface through which LLM agents call external tools.  MCP
gateways sit between the model and the tools, offering a natural enforcement
point for pre-execution policy checks.  Yet published MCP deployments focus
on capability integration (web search, code execution, database access) rather
than safety enforcement.

\textbf{Research question.}  \textit{Do deterministic, pre-execution MCP
guardrails reduce invalid-decision rates for LLM power-pathway planning
agents, and what is the cost in tokens and latency?}

We make three contributions:
(1)~a 22-case benchmark (PowerPathway-22) with five failure modes
and deterministic validator tools;
(2)~an MCP gateway that blocks final submission until prerequisite validators
have been invoked; and
(3)~a live evaluation across 88 runs (4 agents $\times$ 22 cases) with
Gemini~2.5~Flash.  Tool availability reduces the invalid-decision rate
from 36.4\% to 9.1\%, and mandatory enforcement further reduces it to
4.5\%.


% ============================================================
\section{Related Work}
% ============================================================

\subsection{LLM Agents for Engineering}

LLM agents have been applied to robotic design~\cite{makatura2023llm},
chemical synthesis~\cite{boiko2023emergent}, and energy-system
analysis~\cite{zheng2024energy}.  These works highlight the ability of
LLMs to reason qualitatively about engineering trade-offs but do not
address the enforcement of hard numerical constraints.  Schick et
al.~\cite{schick2023toolformer} introduced Toolformer, showing that LLMs
can learn to call external APIs to improve factual accuracy; our work
extends this paradigm by imposing \emph{mandatory sequencing constraints}
on which tools must be called before a final answer may be submitted.
ReAct~\cite{yao2023react} introduced the paradigm of interleaving reasoning
and tool-calling in LLM agents; our G1 agent follows this paradigm but adds
a stateful enforcement layer that restricts which tools may be called and
in what order.

\subsection{Guardrails and Alignment}

Constitutional AI~\cite{bai2022constitutional} and RLHF train models to
avoid undesirable outputs, but provide no formal guarantees at inference
time.  Output-filtering guardrails (e.g., NeMo-Guardrails~\cite{rebedea2023nemo},
Guardrails~AI~\cite{guardrailsai}) post-process model text but cannot
prevent the model from generating a hallucinated numerical answer that
later slips through a lenient regex.  Our approach is complementary:
rather than filtering outputs, we enforce \emph{preconditions on tool calls},
preventing the agent from finalizing a recommendation before deterministic
validators have cleared it.

Doshi et al.~\cite{doshi2026} argue that safety properties of
tool-using agents should be enforced structurally at the tool boundary
rather than through prompt discipline; our MCP gateway instantiates
this principle.  Yang et al.~\cite{yang2025} show that prompts
systematically underspecify intended behavior, motivating enforcement
through typed tool permissions.  Greshake et al.~\cite{greshake2023indirect}
demonstrate that indirect prompt injection can hijack agent workflows;
we observe a related non-adversarial bypass where the model ignores
tools entirely (Section~\ref{sec:validators}).

\subsection{Model Context Protocol}

MCP~\cite{anthropic2024mcp} standardizes tool registration, schema
validation, and result handling for LLM agents.  Existing MCP servers
expose capabilities (databases, web search, code sandboxes) but do not
implement dependency or ordering enforcement.  We extend the gateway
pattern with a stateful \emph{session} that tracks which validators have
been called and blocks downstream calls until prerequisites are satisfied.

\subsection{Power-Pathway Planning}

Grid interconnect capacity, dispatch factors, and capital-cost models are
well-characterized in the power-systems literature~\cite{nrel2023atb}.
However, agent-based tools for power-pathway recommendation are nascent,
and no prior work has evaluated LLM agents against a systematic constraint
benchmark for this domain.


% ============================================================
\section{Method}
% ============================================================

\subsection{Benchmark: PowerPathway-22}

We construct 22 cases representing data-center or industrial facility
planning scenarios with five failure modes (Table~\ref{tab:benchmark}).
Each case specifies load phases (MW), confirmed grid capacity, budget
ceiling, reliability target (\% uptime), and carbon limit (kg\,CO\textsubscript{2}/MWh).
Technology parameters (capacity factor, capex/kW, carbon intensity,
deployment time) are drawn from the NREL Annual Technology Baseline and
the Power Pathway Optimizer prototype~\cite{ppopt}.  All cases and their
ground-truth feasibility labels were constructed by the authors; the
benchmark therefore primarily serves as a controlled testbed for the
guardrail mechanism rather than a representative sample of real-world planning
scenarios.

\begin{table}[t]
\centering
\caption{PowerPathway-22 Benchmark Cases}
\label{tab:benchmark}
\begin{tabular}{lrp{3.8cm}}
\toprule
Failure Mode & Count & Representative challenge \\
\midrule
FEASIBLE           & 6  & Positive control; valid mix exists \\
INFEASIBLE\_PHASE  & 5  & Load exceeds confirmed supply in some phase \\
BUDGET\_VIOLATION  & 4  & Cheapest valid mix exceeds budget ceiling \\
CARBON\_VIOLATION  & 3  & No allowed tech meets carbon target \\
HALLUCINATION\_TRAP& 4  & Missing/ambiguous data tempts false assumption \\
\midrule
Total              & 22 & \\
\bottomrule
\end{tabular}
\end{table}

\subsection{Deterministic Validator Tools}

We implement each hard constraint as a deterministic validator tool.
These validators serve two roles: (1)~as \emph{result invariants},
re-run post-hoc on every agent's final supply mix to compute the
invalid-decision metric; and (2)~as \emph{runtime tools} exposed
to B3 and G1 during planning, where the difference is whether calling
them is optional (B3) or enforced (G1).

The four validators check: \textbf{capacity}, firm (dispatchable)
power $\sum_t \text{CF}_t \times \text{MW}_t \geq \text{peak\_load}$,
with grid capped at the confirmed interconnect; \textbf{cost}, total
capex $\leq$ budget, using NREL ATB~\cite{nrel2023atb} cost rates;
\textbf{carbon}, generation-weighted intensity $\leq$ target; and
\textbf{reliability}, blended availability $\geq$ uptime target, with
a backup bonus for fast-dispatch sources.  For example, a 100\,MW
solar-only proposal at 20\% capacity factor contributes only 20\,MW of
firm capacity, so it fails the capacity validator for a 100\,MW peak-load
case even though its nameplate capacity appears sufficient.  A fifth tool,
\textbf{score\_pathway}, computes a composite score from 0 to 100.
All tools are pure functions with no side effects.

\subsection{MCP Gateway}

The MCP gateway (Fig.~\ref{fig:gateway}) maintains a per-run session object
that records which validators have been invoked.  A blocking rule prevents
\texttt{submit\_final\_recommendation} from executing unless all four
validators have been called.  Attempting to bypass this check returns an
\texttt{ERROR} status with an explanatory message, prompting the agent to
complete the missing validators before resubmitting.

\begin{figure}[t]
\centering
\begin{lstlisting}
# Simplified gateway logic
def call(tool_name, **kwargs):
    if tool_name not in TOOL_REGISTRY:
        return ToolResult("ERROR", "Unknown tool")
    if missing_args(tool_name, kwargs):
        return ToolResult("ERROR", "Missing required args")
    for condition, msg in BLOCKING_RULES[tool_name]:
        if condition(kwargs, session):
            return ToolResult("BLOCKED", msg)
    result = TOOL_REGISTRY[tool_name](**kwargs)
    session.record(tool_name, result)
    return result
\end{lstlisting}
\caption{Simplified MCP gateway enforcement logic. Schema errors return
\texttt{ERROR}; dependency violations return \texttt{BLOCKED}.}
\label{fig:gateway}
\end{figure}

\subsection{Agent Conditions}

We compare four conditions that progressively add structure.  B1 and B2
have no access to validator tools; B3 and G1 both have access to the
same tools but differ in whether usage is enforced.

\textbf{B1 (prompt-only)} issues a single inference call with a system
prompt describing the four constraint types and the required JSON output
format.  No tools are available.

\textbf{B2 (self-critique)} follows B1's initial response with up to two
rounds of structured critique prompts that explicitly enumerate each
constraint, asking the model to check and revise.

\textbf{B3 (tools, no enforcement)} exposes the same six tools as G1
via the model's native function-calling interface, and the system prompt
instructs the agent to call all four validators before recommending.
However, no gateway enforcement exists: the agent \emph{may} submit a
final recommendation without having called any validator.  This condition
isolates the effect of tool availability from the effect of mandatory
enforcement.

\textbf{G1 (MCP-guarded)} is identical to B3 except that all tool calls
pass through the MCP gateway, which blocks
\texttt{submit\_final\_recommendation} unless all four validators have
been called.  This condition tests whether structural enforcement
provides additional safety beyond tool availability with prompting.

\subsection{Evaluation Metrics}

We evaluate each of the 88 runs (4 agents $\times$ 22 cases):
\begin{itemize}
  \item \textbf{invalid\_decision}: the agent's final decision is invalid if
    it either submits a supply mix that violates $\geq$1 hard constraint
    (capacity, cost, carbon, or reliability), as determined by re-running
    all four deterministic validators, or incorrectly declares a feasible
    case infeasible/unknown.  For G1, the gateway structurally prevents
    submission before validators complete, so remaining invalid decisions
    measure whether the \emph{agent correctly interprets validator feedback
    and identifies feasibility}, not merely whether it called the tools.
  \item \textbf{correct\_feasibility}: the agent correctly identified whether
    the case is feasible or infeasible (or unknown for HT-002).  This is the
    primary metric of \emph{reasoning quality}.
  \item \textbf{validators\_called} (B3 and G1): number of the four required
    validators invoked per run.
  \item \textbf{total\_tokens} and \textbf{latency\_sec}: efficiency.
\end{itemize}


% ============================================================
\section{Results}
% ============================================================

\subsection{Overall Safety Metrics}

Table~\ref{tab:main} summarizes the live evaluation results across all
88 runs (4 agents $\times$ 22 cases, Gemini~2.5~Flash).
G1 reduces the invalid-decision rate to 4.5\% (1/22), compared to
36.4\% (8/22) for both B1 and B2.
Notably, self-critique (B2) provides \emph{no improvement} over the
prompt-only baseline: both produce 8 invalid decisions across
identical failure modes.  B3, which has access to the same tools as G1
but without gateway enforcement, achieves 9.1\% (2/22), a substantial
improvement over B1/B2 but still twice G1's rate.  The comparison indicates
that tool availability helps, while enforcement may provide an additional
safety margin.  The \emph{correct feasibility} metric shows a similar
ordering: B1 59.1\%, B2 61.9\%, B3 85.7\%, G1 90.0\%.  Token overhead
for B3 and G1 is 21.8k and 26.1k tokens respectively, reflecting
multi-round function-calling, but latency remains modest (41\,s).

\begin{table}[t]
\centering
\caption{Agent Comparison Across 22 Cases (Live, Gemini~2.5~Flash)}
\label{tab:main}
\begin{tabular}{lrrrrr}
\toprule
Agent & Invalid\% & Correct\% & Avg Tok & Avg Lat (s) & Validators \\
\midrule
B1    & 36.4   &  59.1     &  1{,}161  &  34.5     & N/A  \\
B2    & 36.4   &  61.9     &  8{,}589  &  87.6     & N/A  \\
B3    &  9.1   &  85.7     & 21{,}783  &  41.1     & 3.86 \\
G1    &  4.5   &  90.0     & 26{,}067  &  41.7     & 3.82 \\
\bottomrule
\end{tabular}
\vspace{0.5mm}
\footnotesize Invalid\% denotes invalid final decisions, including both
constraint-violating mixes and incorrect infeasibility/unknown decisions.
Correct\% excludes runs with no parseable feasibility decision
(denominators: B1=22, B2=21, B3=21, G1=20).
\end{table}

\subsection{Breakdown by Failure Mode}

Table~\ref{tab:modes} shows invalid-decision counts by failure mode.
The most striking pattern is the \textsc{FEASIBLE} regime: B1 and B2 each
produce 5/6 invalid decisions: they claim infeasibility for cases
where a valid supply mix exists, unable to find one without validator
guidance.  B3 reduces this to 2/6 and G1 to 1/6, confirming that
validator tools help the model discover feasible mixes that it cannot
find through reasoning alone.  All four agents handle
\textsc{BUDGET\_VIOLATION} and \textsc{CARBON\_VIOLATION} correctly
(0~invalids), because these cases are infeasible and the models
correctly decline to propose a mix.  On \textsc{HALLUCINATION\_TRAP}
cases, B1 and B2 each produce 2/4 invalid decisions; both B3 and
G1 produce 0/4, as validator feedback exposes the missing or
inconsistent data before submission.

\begin{table}[t]
\centering
\caption{Invalid Decisions by Failure Mode (count/total, live)}
\label{tab:modes}
\begin{tabular}{lrrrr}
\toprule
Failure Mode & B1 & B2 & B3 & G1 \\
\midrule
FEASIBLE            & 5/6  & 5/6  & 2/6  & 1/6 \\
INFEASIBLE\_PHASE   & 1/5  & 1/5  & 0/5  & 0/5 \\
BUDGET\_VIOLATION   & 0/4  & 0/4  & 0/4  & 0/4 \\
CARBON\_VIOLATION   & 0/3  & 0/3  & 0/3  & 0/3 \\
HALLUCINATION\_TRAP & 2/4  & 2/4  & 0/4  & 0/4 \\
\midrule
Total               & 8/22 & 8/22 & 2/22 & 1/22 \\
\bottomrule
\end{tabular}
\end{table}

\subsection{Tool Usage Patterns in B3 and G1}
\label{sec:validators}

Both B3 and G1 called validators at similar rates: B3 averaged 3.86
validators per run and G1 averaged 3.82.  When tools are available and
the prompt instructs their use, the model generally
complies, even without enforcement.  The difference between B3 and G1
is therefore not in \emph{whether} validators are called, but in how
the model responds to validator feedback.  B3's 2 invalid
decisions both occurred in \textsc{FEASIBLE} cases despite
calling all four validators: in one case (F-001) the model iterated
repeatedly but could not converge on a capacity-satisfying mix; in
another (F-003) it incorrectly concluded the case was infeasible.
G1 produced one similar error, suggesting that enforcement helps the
model iterate more carefully but does not fully eliminate reasoning
failures.

In one G1 run, the model returned a plain-text response without invoking
any tools; this is a known limitation addressable at the agent level by
requiring tool-call output format.

% ============================================================
\section{Discussion and Threats to Validity}
% ============================================================

\textbf{Construct validity.}  The benchmark cases and ground-truth labels
were constructed by the same author who implemented the guardrails, creating
a risk of confirmation bias in both case design and metric definition.
Independent validation of the benchmark, for example by domain experts or
via crowdsourced annotation, is needed to establish external validity.

\textbf{Internal validity.}  The benchmark spans five failure modes but
remains small (22 cases).  It may not cover all failure patterns encountered
in real deployments.  The deterministic validators are simplified models
that do not capture real-world grid dynamics (e.g., transmission constraints,
curtailment, market pricing).

\textbf{External validity.}  All results use Gemini~2.5~Flash at default
temperature.  Behavior may differ for other model families (GPT-4o,
Claude~3.5, Llama~3) or under non-default temperature settings.  Real
agents face noisier, less structured case descriptions than our benchmark
provides.

\textbf{Cost and benefit.}  B3 and G1 use 21.8k and 26.1k tokens per run
(18.8$\times$ and 22.5$\times$ B1), primarily from multi-round function
calls and model thinking.  However, latency for both is ${\sim}$41\,s, only
1.2$\times$ B1 and actually faster than B2 (87.6\,s), since sequential
critique rounds are expensive.  For 100+\,MW infrastructure decisions
with capital expenditure in the hundreds of millions, the token overhead
is justified given the 4 to 8$\times$ reduction in invalid decisions.

\textbf{Alternative designs.}  A lighter-weight alternative to gateway
enforcement is post-hoc repair: run the agent without safeguards, re-run
the validators on its proposed mix, and feed any violations back as
error messages for the agent to revise on a subsequent turn.  This
avoids MCP-level modification and may match G1 on benchmarks where
validators are pure functions and a wrong intermediate answer is
recoverable.  Pre-execution enforcement remains preferable when tool
calls carry side effects (e.g., procurement, dispatch commitments) or
when the agent's first emitted recommendation is itself acted upon,
which is the regime our gateway targets.  We therefore do not claim that
pre-execution enforcement dominates post-hoc repair; our comparison isolates
enforcement against tool availability without mandatory enforcement.  A
head-to-head comparison of pre-execution enforcement versus post-hoc repair
under matched token budgets is left to future work.

\textbf{Generalizability.}  The core mechanism, a stateful gateway enforcing
prerequisite tool calls, should generalize to other safety-critical agentic
domains (structural design, chemical process design, medical treatment
planning).  Future work should evaluate whether the architecture transfers
without domain-specific tuning.

% ============================================================
\section{Conclusion}
% ============================================================

This paper presents a deterministic MCP guardrail framework for LLM-based
power-pathway planning agents, evaluated in a live experiment with
Gemini~2.5~Flash across 88 runs (4 agents $\times$ 22 cases).  The
primary finding is that providing validator tools (B3) reduces the
invalid-decision rate from 36.4\% to 9.1\%, and adding mandatory
pre-execution enforcement (G1) further reduces it to 4.5\%.
Self-critique (B2) provides no benefit over the prompt-only baseline.
Tool availability accounts for most of the observed improvement, while
enforcement can provide an additional safety margin for constraint-heavy
engineering domains.  Future work will extend the benchmark, evaluate
across multiple model families, and investigate adversarial attacks
targeting gateway bypass.

\bibliographystyle{IEEEtran}
\bibliography{refs}

\end{document}
